%==============================================================================
%  Tugas 3 --- Rangkuman Rencana Riset Doktoral
%  Topik: Online Federated NIDS + Adversarial Training (cross-dataset, cloud)
%==============================================================================
\documentclass[11pt,a4paper]{article}

\usepackage[utf8]{inputenc}
\usepackage[T1]{fontenc}
\usepackage[bahasa]{babel}
\usepackage{geometry}
\geometry{margin=2.5cm}
\usepackage{amsmath,amssymb}
\usepackage{booktabs}
\usepackage{longtable}
\usepackage{array}
\usepackage{enumitem}
\usepackage{ragged2e}
\usepackage{xcolor}
\usepackage{tikz}
\usetikzlibrary{arrows.meta,positioning,calc}
\usepackage{hyperref}
\hypersetup{colorlinks=true,linkcolor=blue,urlcolor=blue,citecolor=blue}

\tikzset{
  spine/.style={-{Stealth[length=3mm]},very thick},
  bone/.style={-{Stealth[length=2mm]},thick,gray!60!black},
  head/.style={rectangle,rounded corners,draw=black,fill=blue!15,
               align=center,font=\small\bfseries,minimum height=1.3cm,text width=2.8cm},
  cat/.style={rectangle,rounded corners,draw=black,fill=green!12,
              align=center,font=\footnotesize\bfseries,text width=2.6cm},
  twig/.style={font=\scriptsize,align=left,text width=3.4cm},
}

\title{\textbf{Rangkuman Rencana Riset Doktoral}\\[4pt]
\large Arsitektur \emph{Online Federated} NIDS dengan \emph{Adversarial Training}\\
untuk Mitigasi \emph{Evasion Attack} dan Adaptasi \emph{Cross-Dataset}\\
pada Lingkungan \emph{Cloud}}
\author{}
\date{}

\begin{document}
\maketitle

%------------------------------------------------------------------
\begin{center}
\begin{tabular}{@{}ll@{}}
\textbf{Nama} & : Hero Yudo Martono \\
\textbf{NRP}  & : 1026900012 \\
\end{tabular}
\end{center}
\vspace{0.4em}
\hrule
\vspace{0.8em}

\noindent\textbf{Judul Topik Riset Doktoral:}\\[2pt]
\emph{``Arsitektur Online Federated NIDS dengan Adversarial Training untuk
Mitigasi Evasion Attack dan Adaptasi Cross-Dataset pada Lingkungan Cloud.''}

\vspace{0.6em}
\noindent\textbf{Benang merah penelitian.} Disertasi ini dibangun bertahap
melalui enam studi yang saling menyambung: (i) deteksi serangan SSH tunggal
sebagai titik awal, (ii) antisipasi ransomware, (iii) ketahanan terhadap
serangan adversarial (\emph{evasion}), (iv) interoperabilitas lintas-dataset
melalui \emph{Semantic Feature Mapping} (SFM), (v) gabungan SFM dengan
\emph{adversarial training}, (vi) adaptasi \emph{online} terhadap
\emph{concept drift}, hingga (vii) desentralisasi melalui \emph{federated
learning}. Semua komponen dimuarakan ke satu arsitektur NIDS yang
\emph{online}, \emph{federated}, tahan \emph{evasion}, dan mampu beradaptasi
lintas-dataset di lingkungan cloud nyata.

%==================================================================
\section{Rumusan Masalah}
%==================================================================

NIDS berbasis \emph{machine learning} menghadapi sejumlah keterbatasan yang
saling berkaitan. Disertasi ini merumuskan masalah berikut:

\begin{enumerate}[label=\textbf{RM\arabic*.},leftmargin=3.2em]
  \item \textbf{Deteksi serangan spesifik.} Bagaimana membangun pendeteksi
  serangan terarah (mis.\ \emph{brute-force} SSH) yang akurat dan dapat
  dioperasikan pada trafik nyata di cloud? [1]

  \item \textbf{Antisipasi ransomware.} Bagaimana memodelkan dan mendeteksi
  perilaku ransomware (mis.\ enkripsi/korupsi berkas) secara dini agar dampak
  dapat dimitigasi? [2]

  \item \textbf{Kerentanan terhadap evasion.} Model NIDS mudah dikelabui
  \emph{adversarial examples}. Bagaimana meningkatkan ketahanan model terhadap
  serangan \emph{evasion} yang \emph{functionality-preserving}? [3,4,5]

  \item \textbf{Generalisasi lintas-dataset.} Skema fitur antar-dataset/tool
  ekstraksi berbeda (CICFlowMeter vs Argus/Bro), sehingga model tidak dapat
  dipindah antar-jaringan. Bagaimana \emph{Semantic Feature Mapping} (SFM)
  menjembatani heterogenitas skema fitur agar model tergeneralisasi
  lintas-dataset? [6,7,8]

  \item \textbf{Ketahanan sekaligus generalisasi.} Dapatkah SFM dan
  \emph{adversarial training} dipadukan sehingga model tahan \emph{evasion}
  sekaligus tetap mampu beradaptasi lintas-dataset? [3,6]

  \item \textbf{Adaptasi online terhadap drift.} Trafik jaringan berubah
  seiring waktu (\emph{concept drift}). Bagaimana model diperbarui secara
  \emph{online} agar tetap akurat tanpa melatih ulang dari nol? [9,10]

  \item \textbf{Desentralisasi \& privasi.} Data lalu lintas antar-organisasi
  tidak boleh dipusatkan. Bagaimana \emph{federated learning} memungkinkan
  pelatihan kolaboratif lintas-klien berskema-fitur-beda tanpa berbagi data
  mentah, dan apa batas kinerjanya di bawah heterogenitas lintas-dataset?
  [11,12,13]

  \item \textbf{Integrasi arsitektur.} Bagaimana menyatukan deteksi,
  ketahanan adversarial, adaptasi online, dan federasi ke dalam satu arsitektur
  NIDS yang utuh dan teruji di lingkungan cloud nyata? [1--13]
\end{enumerate}

%==================================================================
\section{Perbandingan Studi Literatur}
%==================================================================

Tabel~\ref{tab:lit} merangkum studi literatur yang menjadi landasan tiap
komponen disertasi, dikelompokkan menurut tema penelitian.

\renewcommand{\arraystretch}{1.3}
\begin{longtable}{@{}>{\RaggedRight\arraybackslash}p{3.3cm}
                     >{\RaggedRight\arraybackslash}p{8.4cm}
                     >{\RaggedRight\arraybackslash}p{1.9cm}@{}}
\caption{Perbandingan studi literatur per tema penelitian.}
\label{tab:lit}\\
\toprule
\textbf{Studi Literatur (Tema)} & \textbf{Rangkuman} & \textbf{Referensi} \\
\midrule
\endfirsthead
\multicolumn{3}{c}{\tablename\ \thetable\ -- \textit{lanjutan}}\\
\toprule
\textbf{Studi Literatur (Tema)} & \textbf{Rangkuman} & \textbf{Referensi} \\
\midrule
\endhead
\midrule
\multicolumn{3}{r}{\textit{bersambung ke halaman berikutnya}}\\
\endfoot
\bottomrule
\endlastfoot

Deteksi serangan SSH (titik awal) &
Deteksi \emph{brute-force}/SSH berbasis fitur alir pada dataset NIDS modern;
titik awal pipeline deteksi pada trafik cloud nyata. Menggunakan dataset
benchmark NIDS sebagai acuan karakterisasi serangan. &
[1,7,8] \\

Antisipasi ransomware &
Pemodelan perilaku ransomware (enkripsi/overwrite/korupsi metadata) dan deteksi
dini; evaluasi dampak serta pemulihan (RTO) pada simulasi cloud. &
[2] \\

Serangan \& ketahanan adversarial &
Landasan \emph{adversarial examples} (FGSM, PGD, C\&W) dan pemodelan serangan
realistis terhadap NIDS; \emph{adversarial training} sebagai pertahanan;
tinjauan dualitas serangan--pertahanan pada NIDS. &
[3,4,5] \\

Generalisasi \& \emph{domain adaptation} lintas-dataset &
Studi generalisasi lintas-dataset NIDS, standardisasi himpunan fitur
(NetFlow), serta teori adaptasi domain (CORAL, Ben-David) yang melandasi SFM. &
[6,7,8] \\

Semantic Feature Mapping (SFM) + few-shot (fondasi sendiri) &
Pemetaan fitur semantik tervalidasi statistik antar-dataset + kalibrasi
\emph{few-shot}; divalidasi pada trafik cloud nyata (Paper 1 --- fondasi
disertasi). &
[6] \\

Adaptasi \emph{online} / \emph{concept drift} &
Deteksi dan adaptasi terhadap \emph{concept drift} secara \emph{streaming}
(ADWIN/DDM, pembaruan inkremental) agar model tetap mutakhir tanpa latih-ulang
penuh. &
[9,10] \\

\emph{Federated learning} untuk NIDS (survei \& arsitektur) &
Survei konsep, arsitektur, dan strategi agregasi FL-NIDS; FedAvg/FedProx
sebagai dasar; peluang dan tantangan desentralisasi. &
[11,12,14] \\

FL non-IID \& heterogenitas lintas-klien &
Batas FL pada data \emph{non-IID}/heterogen; personalisasi dan penyelarasan
lintas-domain --- relevan dengan temuan kolaps FedAvg pada skema CIC vs UNSW. &
[13,15,16] \\

FL + adversarial robustness &
Ketahanan FL terhadap \emph{adversarial}/poisoning (agregasi robust, SFAT);
menautkan sumbu adversarial dengan sumbu federated. &
[17,18] \\

FL anomaly-based (SOTA pembanding) &
Fed-ANIDS: FL + \emph{autoencoder} (AE/VAE/USAD) berbasis \emph{reconstruction
error}; mencapai kinerja mendekati terpusat pada data homogen (satu-sumber). &
[19] \\

\end{longtable}

\noindent\footnotesize\textit{Catatan kejujuran akademik:} nomor referensi
untuk tema SSH tunggal [1] dan ransomware [2] merujuk naskah/karya penulis
sendiri yang statusnya perlu dilengkapi (judul, venue, tahun) saat finalisasi.
Entri lintas-dataset, adversarial, online, dan federated mengacu pada
literatur terverifikasi (lihat Daftar Referensi). Beberapa entri FL berstatus
pracetak (arXiv) dan DOI-nya perlu diverifikasi ulang sebelum publikasi.

%==================================================================
\section{Hasil Identifikasi Celah (Gap) dan Solusi yang Diusulkan}
%==================================================================

Dari tinjauan literatur (Tabel~\ref{tab:lit}), teridentifikasi lima celah utama.
Tiap celah dipasangkan dengan solusi yang diusulkan disertasi ini.

\renewcommand{\arraystretch}{1.3}
\begin{longtable}{@{}>{\RaggedRight\arraybackslash}p{0.6cm}
                     >{\RaggedRight\arraybackslash}p{6.0cm}
                     >{\RaggedRight\arraybackslash}p{7.0cm}@{}}
\caption{Celah penelitian dan solusi yang diusulkan.}\label{tab:gap}\\
\toprule
\textbf{\#} & \textbf{Celah (Gap) yang Teridentifikasi} & \textbf{Solusi yang Diusulkan} \\
\midrule
\endfirsthead
\toprule
\textbf{\#} & \textbf{Celah (Gap)} & \textbf{Solusi} \\
\midrule
\endhead
\bottomrule
\endlastfoot

G1 & Model NIDS terlatih pada satu dataset tidak tergeneralisasi
ke jaringan lain karena \textbf{skema fitur antar-tool ekstraksi berbeda};
literatur hanya memakai penyelarasan fitur implisit, bukan peta semantik
tervalidasi. &
\textbf{Semantic Feature Mapping (SFM)}: peta fitur semantik tervalidasi
statistik + kalibrasi \emph{few-shot}, memungkinkan transfer model
lintas-dataset tanpa melatih ulang penuh. \\

G2 & NIDS berbasis ML \textbf{rapuh terhadap \emph{evasion}}
(\emph{adversarial examples} yang mempertahankan fungsionalitas serangan);
pertahanan sering diuji dengan serangan tak-realistis. &
\textbf{Adversarial training} dengan serangan \emph{functionality-preserving}
yang realistis; evaluasi ketahanan diukur bersama metrik deteksi normal. \\

G3 & \textbf{SFM dan ketahanan adversarial dikaji terpisah};
belum ada studi apakah interoperabilitas lintas-dataset bertahan di bawah
tekanan adversarial. &
\textbf{SFM + adversarial training terpadu}: menguji apakah ruang fitur
semantik tetap stabil (robust) ketika model dilatih adversarial. \\

G4 & Model statis \textbf{usang akibat \emph{concept drift}};
pelatihan ulang penuh mahal dan lambat untuk operasi \emph{real-time}. &
\textbf{Adaptasi \emph{online}}: deteksi drift + pembaruan model inkremental
pada ruang SFM, menjaga akurasi tanpa latih-ulang dari nol. \\

G5 & FL-NIDS SOTA sukses pada data \textbf{homogen (satu-sumber dipecah)},
tetapi \textbf{FedAvg/FedProx naif divergen} pada heterogenitas lintas-dataset
nyata (prior label \& skema berbeda). &
\textbf{Online Federated NIDS berbasis SFM}: SFM sebagai \emph{enabler}
interoperabilitas + agregasi sadar-heterogenitas; sekaligus memetakan
\emph{batas} FL lintas-dataset (temuan jujur). \\

\end{longtable}

%==================================================================
\section{Konsep Penyelesaian dan Pemetaan Judul Paper}
%==================================================================

Celah di atas diselesaikan melalui rangkaian paper yang saling membangun.
Pemetaan konsep $\rightarrow$ judul paper disajikan pada Tabel~\ref{tab:paper}.

\renewcommand{\arraystretch}{1.3}
\begin{longtable}{@{}>{\RaggedRight\arraybackslash}p{0.9cm}
                     >{\RaggedRight\arraybackslash}p{4.4cm}
                     >{\RaggedRight\arraybackslash}p{7.3cm}
                     >{\RaggedRight\arraybackslash}p{1.0cm}@{}}
\caption{Pemetaan konsep ke judul paper yang direncanakan.}\label{tab:paper}\\
\toprule
\textbf{Paper} & \textbf{Konsep Kunci} & \textbf{Judul (rencana)} & \textbf{Gap} \\
\midrule
\endfirsthead
\toprule
\textbf{Paper} & \textbf{Konsep} & \textbf{Judul} & \textbf{Gap} \\
\midrule
\endhead
\bottomrule
\endlastfoot

P0 & Deteksi SSH + validasi cloud &
Deteksi serangan SSH berbasis ML pada trafik cloud nyata (titik awal pipeline).
& --- \\

P0b & Antisipasi ransomware &
Pemodelan \& deteksi dini perilaku ransomware di lingkungan cloud. & --- \\

P1 & SFM + few-shot &
\emph{Cross-Dataset Generalization of ML-based NIDS via Semantic Feature
Mapping and Few-Shot Adaptation, Validated on Real Cloud Traffic.} & G1 \\

P2 & SFM + adversarial &
\emph{Robust Cross-Dataset NIDS: Adversarial Training on a Semantic Feature
Space.} & G2, G3 \\

P3 & Online / concept drift &
\emph{Online Adaptation of Cross-Dataset NIDS under Concept Drift via
Incremental Learning on the SFM Space.} & G4 \\

P4 & Federated + SFM &
\emph{Federated Cross-Network NIDS via Semantic Feature Mapping: when does
cross-dataset FL work, and why is SFM necessary but not sufficient?} & G5 \\

P5 & Integrasi (disertasi) &
\emph{Online Federated NIDS with Adversarial Training for Evasion Mitigation
and Cross-Dataset Adaptation in Cloud.} & G1--G5 \\

\end{longtable}

%==================================================================
\section{Keunggulan yang Diharapkan (Klaim Kinerja)}
%==================================================================

\noindent\textbf{Indikator kinerja yang diklaim: \emph{Matthews Correlation
Coefficient} (MCC) pada evaluasi \emph{cross-dataset}.} MCC dipilih (bukan
akurasi) karena tahan terhadap ketidakseimbangan kelas yang lazim pada trafik
NIDS, sehingga menjadi ukuran kualitas deteksi yang paling jujur.

\begin{quote}
\textbf{Klaim keunggulan.} \emph{Model NIDS yang dilatih pada dataset sumber
dan dipindah ke dataset target melalui \textbf{Semantic Feature Mapping (SFM)}
mempertahankan MCC cross-dataset yang \textbf{secara signifikan lebih tinggi}
dibandingkan model baseline tanpa SFM (transfer fitur mentah / by-position),
pada tingkat anggaran label target yang sama.}
\end{quote}

\noindent Klaim ini bersifat \emph{terukur} (MCC), \emph{terbatas} (satu
indikator, kondisi cross-dataset), dan \emph{dapat dibantah} (falsifiable):
jika SFM tidak meningkatkan MCC melebihi baseline, klaim gugur.

%==================================================================
\section{Penjelasan Ilmiah atas Klaim (Argumen Deduktif--Nomologis)}
%==================================================================

Klaim keunggulan dijelaskan dengan skema \emph{deductive--nomological} (D--N):
dari hukum/premis umum ($L$) + kondisi awal ($C$), kesimpulan ($E$) menyusul
secara deduktif.

\begin{description}[leftmargin=2.2em,font=\normalfont\itshape]
  \item[Hukum umum $L_1$ (teori adaptasi domain).] Galat model pada domain
  target dibatasi oleh galat sumber ditambah \emph{divergensi distribusi}
  antar-domain (Ben-David dkk.): $\varepsilon_T(h) \le \varepsilon_S(h) +
  d_{\mathcal{H}\Delta\mathcal{H}}(S,T) + \lambda$. Menurunkan divergensi
  menurunkan batas-atas galat target. [8]
  \item[Hukum umum $L_2$ (invariansi semantik).] Jika dua fitur dari dua tool
  ekstraksi mengukur \emph{besaran semantik yang sama}, menyelaraskannya ke satu
  ruang mengurangi perbedaan distribusi \emph{yang bersifat artefak-skema}
  (bukan perbedaan sejati antar-serangan). [6,7]
  \item[Kondisi awal $C_1$.] SFM memetakan fitur CIC dan UNSW ke 9 fitur
  semantik tervalidasi statistik; artefak skema-fitur dihilangkan.
  \item[Kondisi awal $C_2$.] Baseline tanpa SFM menyelaraskan fitur
  \emph{by-position}/mentah, sehingga artefak skema tetap ada.
  \item[Kesimpulan $E$ (deduktif).] Dari $L_1,L_2,C_1,C_2$: ruang SFM memiliki
  $d_{\mathcal{H}\Delta\mathcal{H}}$ lebih kecil daripada ruang baseline
  $\Rightarrow$ batas-atas galat target lebih rendah $\Rightarrow$ MCC
  cross-dataset model SFM $\ge$ MCC baseline. \textbf{Inilah klaim keunggulan.}
\end{description}

\noindent\textbf{Syarat batas (agar argumen tidak \emph{over-claim}).} Argumen
berlaku bila (i) pemetaan semantik benar-benar valid (divalidasi statistik per
fitur), dan (ii) perbedaan sejati antar-serangan tidak ikut terhapus. Jika
salah satu syarat gagal, kesimpulan $E$ tidak dijamin --- sejalan dengan temuan
Paper 4 bahwa SFM \emph{perlu tetapi tidak cukup} ketika heterogenitas berasal
dari distribusi label (bukan sekadar skema fitur).

%==================================================================
\section{Rencana Pembuktian Klaim Keunggulan}
%==================================================================

Pembuktian memakai \textbf{kombinasi simulasi dan eksperimen}, berlapis dari
terkendali ke realistis.

\begin{enumerate}[label=\textbf{\arabic*.},leftmargin=2.2em]
  \item \textbf{Eksperimen offline terkontrol (inti klaim).}
  \begin{itemize}[leftmargin=1.3em]
    \item Data: CSE-CIC-IDS2018 $\leftrightarrow$ UNSW-NB15 (dua arah).
    \item Perlakuan: (a) SFM, (b) baseline tanpa SFM (by-position/mentah).
    \item Protokol: latih di sumber, uji di target; variasikan anggaran label
    target \emph{few-shot} (0, $k$, \dots). Metrik utama \textbf{MCC} +
    recall/precision/F1/balanced-accuracy.
    \item Statistik: 5 \emph{seed} \{13,42,101,202,303\}, rerata$\pm$std,
    \textbf{CI95}, uji signifikansi (paired-$t$ \& Wilcoxon). Klaim didukung
    bila $\Delta$MCC(SFM $-$ baseline) $>0$ signifikan.
  \end{itemize}

  \item \textbf{Eksperimen ketahanan adversarial.} Serangan
  \emph{functionality-preserving} (mis.\ PGD terbatasi pada fitur yang dapat
  dimanipulasi); ukur penurunan MCC dengan vs tanpa \emph{adversarial training}
  pada ruang SFM.

  \item \textbf{Eksperimen online (concept drift).} Aliran data terurut-waktu;
  bandingkan MCC model statis vs model yang diperbarui inkremental saat drift
  terdeteksi.

  \item \textbf{Simulasi \emph{federated}.} Simulasi FedAvg/FedProx \emph{in
  memory} (logika identik node EC2) untuk prediksi hasil sebelum \emph{deploy};
  uji urutan centralized $\ge$ federated $\ge$ local-only.

  \item \textbf{Validasi \emph{deployment} cloud nyata.} AWS multi-EC2
  (aggregator + klien, IP publik + SG ketat, pertukaran bobot via S3); ukur MCC
  pada trafik cloud nyata + biaya komunikasi (rounds, bytes/round) --- menjembatani
  simulasi dengan dunia nyata.
\end{enumerate}

\noindent\textit{Prinsip kejujuran:} semua angka dari eksperimen nyata; hasil
negatif (mis.\ FL kolaps pada heterogenitas lintas-dataset) dilaporkan apa
adanya sebagai temuan ilmiah, bukan disembunyikan.

%==================================================================
\section{Diagram Alur Penelitian (Fishbone)}
%==================================================================

Gambar~\ref{fig:fishbone} menyajikan diagram tulang-ikan (\emph{Ishikawa}):
tulang punggung menuju tujuan akhir (NIDS online-federated yang tahan evasion
dan adaptif lintas-dataset), dengan lima cabang penyebab/komponen yang
masing-masing dialamatkan oleh satu paper.

\begin{figure}[h]
\centering
\begin{tikzpicture}[node distance=4mm]
  % spine
  \node[head] (goal) at (14.5,0)
       {Online Federated NIDS\\ tahan \emph{evasion}\\ \& adaptif\\ cross-dataset};
  \draw[spine] (0,0) -- (goal.west);
  % upper categories
  \node[cat] (c1) at (2.2,2.2) {SFM\\ (interoperabilitas)};
  \node[cat] (c3) at (7.0,2.2) {Online\\ (concept drift)};
  \node[cat] (c5) at (11.5,2.2) {Deteksi awal\\ (SSH, ransomware)};
  % lower categories
  \node[cat] (c2) at (2.2,-2.2) {Adversarial\\ (evasion)};
  \node[cat] (c4) at (7.0,-2.2) {Federated\\ (privasi)};
  % bones from category to spine
  \draw[bone] (c1.south) -- (2.2,0);
  \draw[bone] (c3.south) -- (7.0,0);
  \draw[bone] (c5.south) -- (11.5,0);
  \draw[bone] (c2.north) -- (2.2,0);
  \draw[bone] (c4.north) -- (7.0,0);
  % twigs (sub-items) -- upper
  \node[twig,anchor=west] at (2.6,3.3) {$\cdot$ peta fitur tervalidasi (P1)};
  \node[twig,anchor=west] at (2.6,3.0) {$\cdot$ few-shot calibration};
  \node[twig,anchor=west] at (7.4,3.3) {$\cdot$ deteksi drift (P3)};
  \node[twig,anchor=west] at (7.4,3.0) {$\cdot$ update inkremental};
  \node[twig,anchor=west] at (11.9,3.3) {$\cdot$ brute-force SSH (P0)};
  \node[twig,anchor=west] at (11.9,3.0) {$\cdot$ perilaku ransomware (P0b)};
  % twigs -- lower
  \node[twig,anchor=west] at (2.6,-3.0) {$\cdot$ adversarial training (P2)};
  \node[twig,anchor=west] at (2.6,-3.3) {$\cdot$ serangan realistis (PGD)};
  \node[twig,anchor=west] at (7.4,-3.0) {$\cdot$ FedAvg/FedProx (P4)};
  \node[twig,anchor=west] at (7.4,-3.3) {$\cdot$ validasi AWS multi-EC2};
\end{tikzpicture}
\caption{Diagram \emph{fishbone} (Ishikawa) alur penelitian doktoral: lima
komponen (SFM, adversarial, online, federated, deteksi awal) menyatu menuju
arsitektur NIDS \emph{online-federated} yang tahan \emph{evasion} dan adaptif
lintas-dataset. Label P0--P5 menautkan tiap cabang ke paper pada
Tabel~\ref{tab:paper}.}
\label{fig:fishbone}
\end{figure}

%==================================================================
\section*{9. Daftar Referensi}
\addcontentsline{toc}{section}{Daftar Referensi}
%==================================================================
\begin{enumerate}[label={[\arabic*]},leftmargin=2.6em]

\item H.~Y.~Martono \emph{et al.}, ``Deteksi serangan SSH berbasis
\emph{machine learning} pada trafik cloud nyata,'' naskah kerja penulis
(Paper SSH), \emph{status terbit dilengkapi saat finalisasi}.

\item H.~Y.~Martono \emph{et al.}, ``Pemodelan dan deteksi dini perilaku
ransomware pada lingkungan cloud,'' naskah kerja penulis (Paper Ransomware),
\emph{status terbit dilengkapi saat finalisasi}.

\item I.~J.~Goodfellow, J.~Shlens, C.~Szegedy, ``Explaining and harnessing
adversarial examples,'' in Proc.\ ICLR, 2015.

\item A.~Madry, A.~Makelov, L.~Schmidt, D.~Tsipras, A.~Vladu, ``Towards deep
learning models resistant to adversarial attacks,'' in Proc.\ ICLR, 2018.

\item G.~Apruzzese, M.~Andreolini, L.~Ferretti, M.~Marchetti, M.~Colajanni,
``Modeling realistic adversarial attacks against network intrusion detection
systems,'' ACM Trans.\ Privacy Security, 2022.

\item H.~Y.~Martono, I.~Syarif, F.~A.~Saputra, ``Cross-Dataset Generalization
of ML-based Network Intrusion Detection via Semantic Feature Mapping and
Few-Shot Adaptation, Validated on Real Cloud Traffic,'' naskah kerja (Paper 1),
\emph{status terbit dilengkapi saat finalisasi}.

\item M.~Sarhan, S.~Layeghy, M.~Portmann, ``Towards a standard feature set for
network intrusion detection system datasets,'' Mobile Netw.\ Appl., vol.~27,
no.~1, pp.~357--370, 2022.

\item S.~Ben-David, J.~Blitzer, K.~Crammer, A.~Kulesza, F.~Pereira,
J.~W.~Vaughan, ``A theory of learning from different domains,'' Mach.\ Learn.,
vol.~79, no.~1--2, pp.~151--175, 2010.

\item S.~Layeghy, M.~Gallagher, M.~Portmann, ``Benchmarking the benchmark ---
Comparing synthetic and real-world network IDS datasets,'' J.\ Inf.\ Secur.\
Appl., vol.~80, art.~103689, 2024.

\item R.~Doriguzzi-Corin, D.~Siracusa, ``FLAD: Adaptive Federated Learning for
DDoS Attack Detection,'' arXiv:2205.06661, 2022.

\item A.~Khraisat, A.~Alazab, T.~Jan, A.~Gopez, ``Survey on Federated Learning
for Intrusion Detection System: Concept, Architectures, Aggregation Strategies,
Challenges, and Future Directions,'' ACM Comput.\ Surv., 2024.
doi:10.1145/3687124.

\item J.~L.~Hern\'andez-Ramos \emph{et al.}, ``Intrusion Detection Based on
Federated Learning: A Systematic Review,'' ACM Comput.\ Surv., 2025.
doi:10.1145/3731596.

\item L.~Lavaur, Y.~Busnel, F.~Autrel, ``Highlighting the Limits of Federated
Learning in Intrusion Detection,'' in Proc.\ IEEE ICDCS, 2024.
doi:10.1109/icdcs60910.2024.00135.

\item A.~A.~Wardana, P.~Sukarno, ``Taxonomy and Survey of Collaborative
Intrusion Detection System using Federated Learning,'' ACM Comput.\ Surv.,
2024. doi:10.1145/3701724.

\item S.~Gurpreet \emph{et al.}, ``Sentinel: Dynamic Knowledge Distillation for
Personalized Federated Intrusion Detection in Heterogeneous IoT Networks,''
arXiv:2510.23019, 2025.

\item M.~Khalid, U.~Zukaib, M.~Al-Rakhami, A.~M.~Alamri, ``FedCross-VAN:
Federated Cross-Domain Behavior Alignment for VANET Intrusion Detection,''
Trans.\ Emerg.\ Telecommun.\ Technol., vol.~37, 2025. doi:10.1002/ett.70309.

\item ``Combating Exacerbated Heterogeneity for Robust Models in Federated
Learning (SFAT),'' arXiv:2303.00250, 2023.

\item E.~M.~Ennaji, S.~E.~Hajla, Y.~Maleh, S.~Mounir, ``Adversarially robust
federated deep learning models for intrusion detection in IoT,'' Indonesian J.\
Electr.\ Eng.\ Comput.\ Sci., vol.~37, no.~2, p.~937, 2024.
doi:10.11591/ijeecs.v37.i2.pp937-947.

\item M.~Janati Idrissi, H.~Alami, A.~El Mahdaouy, A.~El Mekki, S.~Oualil,
Z.~Yartaoui, I.~Berrada, ``Fed-ANIDS: Federated learning for anomaly-based
network intrusion detection systems,'' Expert Syst.\ Appl., vol.~234,
art.~121000, 2023. doi:10.1016/j.eswa.2023.121000.

\end{enumerate}

\end{document}
